\documentclass[11pt,a4paper]{article}
\usepackage[margin=25mm]{geometry}
\usepackage[T1]{fontenc}
\usepackage{lmodern,microtype}
\usepackage{amsmath,amssymb,amsthm,mathtools}
\usepackage{booktabs,tabularx}
\usepackage{graphicx}
\usepackage{pgfplots}
\pgfplotsset{compat=1.18}
\usepackage{hyperref}
\hypersetup{colorlinks=true,linkcolor=blue,citecolor=blue,urlcolor=blue}
\newtheorem{theorem}{Theorem}
\newtheorem{lemma}{Lemma}
\newtheorem{proposition}{Proposition}
\newtheorem{remark}{Remark}
\newtheorem{corollary}{Corollary}
\newcommand{\sat}{\operatorname{sat}}
\newcommand{\sgn}{\operatorname{sgn}}
\newcommand{\R}{\mathbb{R}}
\title{Competing Residence Times and Seven Limit Cycles\\[3pt]\large in a Periodically Forced Three-Zone Equation}
\author{Chaoyu Hu \and Tianrun Yu}
\date{13 October 2026}
\begin{document}
\maketitle
\begin{abstract}
We consider a periodically forced, symmetric, three-zone scalar
piecewise-linear equation with contracting outer zones and an expanding
middle zone. An earlier conjecture proposed that this family admits at most
five limit cycles. We give a rational counterexample and prove a stronger
parameter-uniform statement: every system in an explicitly specified open
cube about $(A,B,\mu)=(9,3/2,181363/100000)$ possesses at least seven
isolated $2\pi$-periodic orbits, including at least four locally attracting
and three locally repelling ones. A residence-time identity and a
square-root law at grazing explain the geometry that led to the example.
The proof uses exact affine flows, strict interval bounds for Poincar\'e
displacements and switching events, scalar order, and half-period
symmetry. We also obtain three distinct critical residence-time profiles.
No sharp global upper bound is claimed.
\end{abstract}
\noindent\textbf{Keywords:} piecewise-linear differential equations; limit cycles; residence time; grazing bifurcations; periodic solutions. \quad
\textbf{MSC 2020:} 34C25, 34A34, 37C27.

\section{Introduction}\label{sec:intro}
We study the periodically forced scalar equation
\begin{equation}\label{eq:main}
 \dot x=-Ax+(A+B)\sat(x)+\mu\sin t,
 \qquad A,B,\mu>0,
 \qquad
 \sat(x)=
 \begin{cases}
   -1,&x<-1,\\
   x,&|x|\le1,\\
   1,&x>1.
 \end{cases}
\end{equation}
Its solutions define trajectories on the cylinder
$(t\bmod 2\pi,x)$. A limit cycle means an isolated periodic
trajectory on this cylinder. Despite the explicit affine flow
in each zone, the number of periodic solutions depends on the
switching times between zones.

Bravo, Carmona, Fern\'andez and Ojeda~\cite{Bravo2025}
proved a parameter-independent finite upper bound and exhibited
parameters with five limit cycles. They proposed five as the
largest possible number. A more general, substantially larger
bound appears in~\cite{HuYuBound}. We give a counterexample
to the proposed bound and prove that the phenomenon persists
on an explicit parameter set.

\begin{theorem}[Seven isolated cycles on an explicit parameter cube]
\label{thm:main}
Let
\[
  (A_0,B_0,\mu_0)
  =\left(9,\frac32,\frac{181363}{100000}\right).
\]
For every $(A,B,\mu)$ satisfying
\begin{equation}\label{eq:param-cube}
  |A-A_0|<10^{-12},\qquad
  |B-B_0|<10^{-12},\qquad
  |\mu-\mu_0|<10^{-12},
\end{equation}
equation~\eqref{eq:main} has at least seven isolated
$2\pi$-periodic trajectories. Among them one can choose at
least four locally attracting and three locally repelling
trajectories. The antisymmetric central trajectory is
hyperbolic repelling; no hyperbolicity is asserted for the
selected noncentral trajectories.
\end{theorem}

In particular, $(A,B,\mu)=(A_0,B_0,\mu_0)$ disproves the
conjectured universal upper bound of five. The theorem is
a \emph{lower} bound for an explicit parameter region:
it does not assert that seven is the maximum, either at
these parameters or globally.

The proof and the search leading to the parameters have
different roles. To locate a candidate, we consider the
mean-parametrized periodic response of a biased equation.
Its derivative is controlled by the time spent in the
expanding middle zone. As the mean changes, the lower
and upper exterior residence times vary in opposite
directions; near grazing one of them vanishes at
square-root order. This produces the additional bending
of the response curve discussed in Sections~\ref{sec:response}
and~\ref{sec:discovery}. The existence proof instead uses
strict Poincar\'e displacement signs, the ordered scalar
flow, and half-period symmetry (Sections~\ref{sec:witness}
and~\ref{sec:proof}). Quantitative control of both grazing
barriers and the central multiplier makes the whole cube
in~\eqref{eq:param-cube} valid.

The supporting computation uses directed decimal interval
arithmetic rather than a formally verified proof assistant.
An independent Arb/FLINT replication would be a valuable
additional check before submission.
 
\section{Residence times and the periodic response}\label{sec:response}
For equation~\eqref{eq:main}, write
\begin{equation}\label{eq:g}
  \dot x+g(x)=\mu\sin t,\qquad
  g(x)=Ax-(A+B)\sat(x).
\end{equation}
The outer regions have $g'=A>0$, hence perturbations decay at rate $A$;
the middle region has $g'=-B<0$, hence they grow at rate $B$.
A single flight therefore tells us nothing decisive about the full-period
response: its total amplification depends on how long it spends in each zone.

\subsection{The residence-time identity}
To vary periodic solutions continuously without restricting ourselves to
actual zeros of the original equation, introduce an auxiliary bias:
\begin{equation}\label{eq:bias}
 \dot x+g(x)=\mu\sin t+\lambda,
 \qquad x(t+2\pi)=x(t).
\end{equation}
The periodic solutions of this auxiliary family admit a global mean
parametrization, as established in \cite[Section~2]{HuYuBound}:
\[
 \xi=\frac1{2\pi}\int_0^{2\pi}x(t)\,dt,
 \qquad \lambda=\Lambda(\xi,\mu),
 \qquad z(t)=\partial_\xi x(t;\xi,\mu)>0.
\]
In this parametrization the original equation is simply the horizontal
section $\Lambda=0$. The half-period odd symmetry implies
\begin{equation}\label{eq:odd}
 \Lambda(-\xi,\mu)=-\Lambda(\xi,\mu),\qquad\Lambda(0,\mu)=0.
\end{equation}
The central orbit is always present; other zeros, if any, occur in
opposite-mean pairs.

Define the middle-zone residence time
\begin{equation}\label{eq:T}
 T(\xi,\mu)=\operatorname{meas}\{t\in[0,2\pi):|x(t;\xi,\mu)|<1\}
\end{equation}
and the critical duration $\rho=2\pi A/(A+B)$.
Differentiating \eqref{eq:bias} with respect to $\xi$ at transverse
crossings gives $\dot z+g'(x)z=\Lambda_\xi$. Because $z>0$, division by $z$
and integration around the periodic orbit yield
\begin{equation}\label{eq:key}
\boxed{\quad
 \Lambda_\xi J=(A+B)(\rho-T),\qquad
 J=\int_0^{2\pi}\frac{dt}{z(t)}>0.
\quad}
\end{equation}
Indeed, the integral of $\dot z/z$ vanishes, while
$\int g'(x)\,dt=A(2\pi-T)-BT$. The identity extends across isolated
grazing contacts by continuity \cite{HuYuBound}.

This identity is the basic reason to study durations rather than explicitly
solve every periodic boundary-value problem: the sign of
$\Lambda_\xi$ depends \emph{only} on whether the residence time is above or
below $\rho$. At $T>\rho$ the response curve descends; at $T<\rho$
it ascends. A five-cycle upper bound would follow from a sufficiently simple
shape of these sign regions. Our example violates that shape.

There is an independent dynamical interpretation. At a genuine periodic
solution of the original equation, the Floquet multiplier of the return map
$P(y)=x(2\pi;0,y)$ is
\begin{equation}\label{eq:multiplier}
\boxed{\quad P'(y)
 =\exp\!\left(-\int_0^{2\pi}g'(x(t))\,dt\right)
 =\exp\bigl((A+B)(T-\rho)\bigr).\quad}
\end{equation}
The same critical duration $T=\rho$ therefore separates contraction from
expansion of the full-period map. It is not necessary to numerically solve a
separate variational equation after the switching durations have been found.

\subsection{The lower and upper visits push in opposite directions}
Suppose a positive-mean trajectory visits both exterior regions. Denote by
$s$ its total time below $-1$ and by $v$ its total time above $1$.
Then identically
\begin{equation}\label{eq:compete}
 T=2\pi-s-v,\qquad T_\xi=-s_\xi-v_\xi.
\end{equation}
Strict order $z>0$ means that raising $\xi$ lifts the entire graph of the
orbit. Thus the time below $-1$ decreases while the time above $1$
increases:
\begin{equation}\label{eq:opposed}
 s_\xi<0,\qquad v_\xi>0
\end{equation}
on a transverse four-crossing itinerary. Consequently $-s_\xi>0$ and
$-v_\xi<0$ compete. No sign follows for their sum. This is the precise
obstruction to a global monotonicity argument for the residence time.

The contribution from the lower exterior visit becomes particularly sharp
as that visit disappears. This can be quantified without a complete
classification of periodic solutions.

\begin{proposition}[Square-root loss of the lower exterior visit]\label{prop:grazing}
Fix $A,B,\mu>0$. Suppose the mean-parametrized periodic family
\eqref{eq:bias} has a nondegenerate lower grazing at $(t_g,\xi_g)$:
\[
 x(t_g;\xi_g,\mu)=-1,\qquad
 x_t(t_g;\xi_g,\mu)=0,\qquad
 \gamma:=x_{tt}(t_g;\xi_g,\mu)=\mu\cos t_g>0.
\]
Assume the upper exterior visit is present with transverse crossings and
this is the only lower grazing in a neighborhood. With
$z_g=\partial_\xi x(t_g;\xi_g,\mu)>0$, the duration of the disappearing
lower visit satisfies
\begin{equation}\label{eq:sqrt}
 \boxed{\quad
 s(\xi)=2\sqrt{\frac{2z_g}{\gamma}}\,
          \sqrt{\xi_g-\xi}+o\bigl(\sqrt{\xi_g-\xi}\bigr)
 \quad}\qquad(\xi\uparrow\xi_g).
\end{equation}
In particular, the middle residence time has a strict local maximum at
the grazing fibre.
\end{proposition}
\begin{proof}
Since the vector field is continuous across $x=-1$, and the velocity
vanishes at the contact, both one-sided second derivatives agree and equal
$\mu\cos t_g=\gamma$. The family is differentiable in its mean at the
contact. For $\delta=\xi_g-\xi\downarrow0$ and
$\tau=t-t_g$, the local expansion is
\[
 x(t_g+\tau;\xi_g-\delta,\mu)+1
   =-z_g\delta+\frac\gamma2\tau^2
     +o(\delta+\tau^2).
\]
The two neighboring intersections with $-1$ therefore occur at
$\tau_\pm=\pm\sqrt{2z_g\delta/\gamma}+o(\sqrt\delta)$, and their
time separation is \eqref{eq:sqrt}.
The upper-zone duration $v(\xi)$ is differentiable near the contact,
with $v_\xi>0$ by strict pointwise order and transverse crossings.
For $\xi<\xi_g$ close to the contact, 
$T(\xi)=2\pi-v(\xi)-s(\xi)$ lies below
$T(\xi_g)=2\pi-v(\xi_g)$ because the square-root loss dominates the
$O(\delta)$ upper-duration change. For $\xi>\xi_g$ the lower visit is absent,
so $T(\xi)=2\pi-v(\xi)<T(\xi_g)$. Hence the grazing fibre is a
strict local maximum of the residence time.
\end{proof}

\begin{remark}
The proposition establishes a \emph{local cusp in $T$ at a generic grazing},
not an automatic additional periodic orbit. Zeros of the original equation
still require $\Lambda=0$, and response extrema require $T=\rho$.
What the cusp supplies is a geometric way for $T$ to turn sharply, and
therefore a reason that multiple crossings of the critical level $\rho$
can occur when the exterior damping is sufficiently asymmetric.
\end{remark}

\section{Locating the seven-cycle parameters}\label{sec:discovery}
The residence-time identity suggests a practical way of locating additional
response folds. The square-root grazing mechanism rules out a simple global
monotonicity argument, and the switching geometry reduces the search to
a finite expression in crossing times and velocities.

\subsection{A four-crossing derivative formula}
Take successive switching times $t_1<t_2<t_3<t_4$, in the order
$+1,-1,-1,+1$, and let
\begin{equation}\label{eq:durations}
 r=t_2-t_1,\quad s=t_3-t_2,\quad
 u=t_4-t_3,\quad v=t_1+2\pi-t_4.
\end{equation}
The intervals alternate between middle, lower, middle, and upper zones.
Here $T=r+u$ and $s+v=2\pi-T$. On a response critical orbit
$T=\rho$, so $u=\rho-r$ and $v=2\pi-\rho-s$: the four durations reduce to
two free variables.

Implicit differentiation of each transverse crossing equation gives
\begin{equation}\label{eq:crossdiff}
 t_{i,\xi}=-\frac{z(t_i)}{x_t(t_i)}.
\end{equation}
Along middle flights $z$ is multiplied by $e^{B\ell}$, whereas along
exterior flights it is multiplied by $e^{-A\ell}$, at each critical fibre
where $\Lambda_\xi=0$. Set $w_i=x_t(t_i)\ne0$, $z_1=z(t_1)>0$.
The identity $T=t_2-t_1+t_4-t_3$ then gives the explicit four-speed test
\begin{equation}\label{eq:H}
\boxed{\quad
 T_\xi=z_1\mathcal H,\qquad
 \mathcal H=\frac1{w_1}-\frac{e^{Br}}{w_2}
 +\frac{e^{Br-As}}{w_3}-\frac{e^{B(r+u)-As}}{w_4}.
 \quad}
\end{equation}
The sign of the residence derivative has now become the sign of a finite
expression in the four crossing velocities. Using the flight matching
relations of \cite[Section~3]{HuYuBound}, admissible critical trajectories
can be parametrized by a compatibility curve in $(r,s)$. The search is thus
for a change of sign on a physically admissible curve, not a blind search
for zeros of the full-period map.

\subsection{Selection of the rational parameters}
Initial numerical tests with weak exterior damping suggested a simple sign
for $\mathcal H$, but that sign failed for stronger damping. A physical
negative value arose with the convenient rational slopes
\[
 A=9,\qquad B=\frac32.
\]
Following the associated response-critical curve numerically then revealed
a shallow \emph{valley} in the required forcing amplitude. Its minimum was
near $\mu_{\min}\simeq1.813584$, while the symmetric critical fibre
occurred at $\mu_c\simeq1.81363456$. The narrow strip
$\mu_{\min}<\mu<\mu_c$ is precisely where a fixed-amplitude level can
intersect the three-zone critical curve twice. We selected the simple
rational value
\begin{equation}\label{eq:mu}
 \mu=\frac{181363}{100000}=1.81363
\end{equation}
in that strip. The two threshold values are exploratory numerical
approximations, not rigorously enclosed bifurcation parameters.

At this amplitude, the numerical response has the following three turning
points on the positive-mean half-line:
\begin{table}[ht]
\centering
\caption{Numerical response extrema; these values motivated the witness but are not proof inputs.}\label{tab:turns}
\begin{tabular}{@{}r r l l@{}}\toprule
Mean $\xi$ & $\Lambda(\xi)$ & Type & Itinerary\\\midrule
$0.00427530$ & $-5.83\cdot10^{-7}$ & local minimum & four crossings\\
$0.02615002$ & $+1.20\cdot10^{-5}$ & local maximum & four crossings\\
$0.04964750$ & $-2.30\cdot10^{-5}$ & local minimum & two crossings\\\bottomrule
\end{tabular}
\end{table}
Beginning from $\Lambda(0)=0$, the graph descends, rises, descends,
and rises again. This is exactly the pattern needed for \emph{three positive
response zeros}, rather than the two allowed by the old five-cycle guess.
The first two extrema belong to the genuine three-zone regime; the third
is the ordinary two-zone minimum after the lower exterior visit vanishes.

\begin{figure}[ht]
\centering
\begin{tikzpicture}
\begin{axis}[width=.94\linewidth,height=6.4cm,grid=major,
 grid style={gray!18},xmin=0,xmax=.08,
 xlabel={Mean $\xi$},ylabel={$10^5\Lambda(\xi,\mu)$},
 scaled x ticks=false,xtick={0,.02,.04,.06,.08},
 label style={font=\small},tick label style={font=\small}]
\addplot[very thick,blue!65!black,mark=none]
 table[x=xi,y=scaled_lambda,col sep=comma]{figures/response_data.csv};
\addplot[dashed,black!70] coordinates {(0,0)(.08,0)};
\addplot[only marks,mark=x,mark size=3pt,very thick,red!75!black]
 coordinates {(.007659607,0) (.037207634,0) (.065677290,0)};
\end{axis}
\end{tikzpicture}
\caption{Exploratory response curve at the rational witness parameters.
Its three positive zeros are verified \emph{independently} by the strict
return-map signs in Table~\ref{tab:signs}. The plotted data are not used
as a mathematical certificate.}\label{fig:response}
\end{figure}

The geometry serves a different purpose from the proof. It tells us
\emph{where} to look and what crossing structure to expect. Near grazing,
however, the first response extremum is tiny, and ordinary time stepping
can give false fixed points. To prove the claim we abandon the plotted
response values and return to the original unbiased equation.

\section{The rational witness and its Poincar\'e map}\label{sec:witness}
Our witness is
\begin{equation}\label{eq:witness}
\boxed{\ \dot x=-9x+\frac{21}{2}\sat(x)+\frac{181363}{100000}\sin t.\ }
\end{equation}
Equivalently, $\dot x+g(x)=\mu\sin t$ with
\begin{equation}\label{eq:gparts}
g(x)=\begin{cases}9x+21/2,&x<-1,\\-3x/2,&|x|\le1,\\9x-21/2,&x>1.\end{cases}
\end{equation}
Note that the vector field is continuous and globally Lipschitz even at $x=\pm1$; existence, uniqueness, and continuous dependence are therefore immediate. For initial value $x(0)=y$, put
\begin{equation}\label{eq:P}
P(y)=x(2\pi;0,y),\qquad D(y)=P(y)-y.
\end{equation}
Fixed points of $P$ are precisely the required periodic solutions.

\begin{proposition}[Certified rational witness]\label{prop:witness}
At $(A,B,\mu)=(A_0,B_0,\mu_0)$, equation~\eqref{eq:main}
has at least seven distinct isolated $2\pi$-periodic
trajectories.
\end{proposition}

For orientation only, an independent high-precision computation gives three initial values above the antisymmetric center and their characteristic multipliers:
\begin{table}[ht]
\centering
\caption{Numerical approximations (not used in the proof).}\label{tab:roots}
\begin{tabular}{@{}r r l@{}}\toprule
$x(0)$ & $P'(x(0))$ & Qualitative stability \\\midrule
$-0.569203954368034$ & $0.9914612791$ & attracting \\
$-0.556760196081272$ & $1.5169801260$ & repelling \\
$-0.555155990304937$ & $0.3330479441$ & attracting \\\bottomrule
\end{tabular}
\end{table}
The second solution approaches $-1$ very closely at its minimum, while the third stays above it and has only two threshold crossings. The mathematical problem is to prove that these are \emph{three genuine and distinct} zeros, not numerical artifacts caused by near-grazing.

\section{Proof of the main theorem}\label{sec:proof}
\subsection{Exact affine flows and switching events}
In a region where $g(x)=cx+d$, the flow from $x(t_0)=y_0$ is
\begin{align}\label{eq:flight}
\Phi_{c,d}(t;t_0,y_0)={}&e^{-c(t-t_0)}y_0-\frac dc\bigl(1-e^{-c(t-t_0)}\bigr)\\
&+\frac{\mu}{1+c^2}\Bigl[c\sin t-\cos t
-e^{-c(t-t_0)}(c\sin t_0-\cos t_0)\Bigr].\notag
\end{align}
Here $c=-3/2$ for the middle zone and $c=9$ for the outer zones. An exact orbit is obtained by joining these flows at the \emph{first} hitting times of $\pm1$. Only the transcendental evaluations and hitting-time brackets need interval enclosures.

There is a particularly simple observation that makes the event audit understandable.
\begin{lemma}[A threshold cannot be crossed repeatedly]\label{lem:cross}
Fix a threshold $\ell\in\{-1,1\}$ and a trajectory of \eqref{eq:witness}. At every transverse encounter with $x=\ell$,
\begin{equation}\label{eq:threshold}
\dot x=\mu\sin t-g(\ell)=\mu\sin t+\frac32\ell.
\end{equation}
For a \emph{periodic} trajectory, the level has at most one upward and one downward transverse crossing per period. For an arbitrary shooting trajectory, any change in the number of level crossings as the initial value varies can occur only at a tangency $\dot x=0$ at that level.
\end{lemma}
\begin{proof}
The times on the phase circle where the right side of \eqref{eq:threshold} is positive form one connected arc (or the empty set); the negative times form the complementary connected arc. Upward and downward level crossings of a continuous scalar trajectory necessarily alternate. On a periodic trajectory the crossings alternate cyclically; two upward crossings would require a downward crossing within the connected positive arc, which is impossible. The same applies to downward crossings. For an arbitrary shooting trajectory, transverse crossings persist under a small variation in the initial state by the implicit function theorem. A change of their number away from the endpoints of the time interval therefore requires a nontransverse encounter. In our shooting brackets, the end states at $t=0$ and $t=2\pi$ remain strictly away from both switching levels, so endpoint loss is also excluded.
\end{proof}

A second useful shape identity holds within each affine zone. If $\dot x+cx+d=\mu\sin t$ and $w(t)=e^{c(t-t_0)}\dot x(t)$, then
\begin{equation}\label{eq:shape}
\boxed{\ w'(t)=\mu e^{c(t-t_0)}\cos t.\ }
\end{equation}
Thus the velocity cannot develop a hidden oscillation inside a flight: it has a known monotonicity pattern on the arcs determined by $\pi/2$ and $3\pi/2$. The two identities also make it possible to check the absence of tangencies across an entire interval of initial values, rather than merely accepting the crossings returned by a numerical event detector.

\subsection{Three disjoint fixed-point brackets}
The six initial values are rational decimal numbers. Table~\ref{tab:signs} gives deliberately widened enclosures for their full-period displacements. The enclosures contain the raw intervals produced by the script \texttt{proof/certify\_seven.py}; no numerical root estimates are used.
\begin{table}[ht]
\centering
\caption{Strict outward bounds for $D(y)=P(y)-y$. Positive/negative endpoints occur in three disjoint pairs.}\label{tab:signs}
\begin{tabular}{@{}r r c@{}}\toprule
Initial value $y$ & Certified enclosure for $D(y)$ & Sign \\\midrule
$-0.571$ & $[1.13334,\ 1.13357]\cdot10^{-5}$ & $+$ \\
$-0.567$ & $[-2.57304,\ -2.57281]\cdot10^{-5}$ & $-$ \\
\addlinespace[2pt]
$-0.55677$ & $[-4.97010,\ -4.96788]\cdot10^{-6}$ & $-$ \\
$-0.55675$ & $[5.38442,\ 5.38665]\cdot10^{-6}$ & $+$ \\
\addlinespace[2pt]
$-0.55517$ & $[9.31073,\ 9.31075]\cdot10^{-6}$ & $+$ \\
$-0.55514$ & $[-1.07074,\ -1.07072]\cdot10^{-5}$ & $-$ \\\bottomrule
\end{tabular}
\end{table}
Since $D$ is continuous, each interval
\begin{equation}\label{eq:intervals}
I_1=[-0.571,-0.567],\quad
I_2=[-0.55677,-0.55675],\quad
I_3=[-0.55517,-0.55514]
\end{equation}
contains at least one fixed point. The three intervals do not overlap.

For the first two pairs the certified zone sequence is
\[
\underbrace{\mathrm{middle}\to\mathrm{lower}\to\mathrm{middle}}_{\text{two $-1$ crossings}}
\to\underbrace{\mathrm{upper}\to\mathrm{middle}}_{\text{two $+1$ crossings}};
\]
for the last pair it is $\mathrm{middle}\to\mathrm{upper}\to\mathrm{middle}$. Each event is enclosed between rational decimal times at which the exact flight lies on opposite sides of the threshold; the crossing speed has a strict sign. The last interval has, additionally, a certified initial minimum strictly \emph{above} $-1$. The shape formula \eqref{eq:shape} rules out an earlier encounter at the endpoint trajectories.

The remaining concern is a possible hidden change of itinerary \emph{between} the rational shooting endpoints. This can now be excluded by a finite argument. A grazing at the level $\ell\in\{-1,1\}$ can occur only at one of the two phases solving
\begin{equation}\label{eq:grazingphases}
\mu\sin t+\tfrac32\ell=0.
\end{equation}
For the chosen $\mu$, the four candidate phases, rounded here for readability, are
\[
0.9738759264,\quad2.1677167272\quad(\ell=-1),\qquad
4.1154685800,\quad5.3093093808\quad(\ell=+1).
\]
The interval checker first certifies each phase by bracketing the root of \eqref{eq:grazingphases}; it then evaluates $x(t;y)-\ell$ at that phase for \emph{both} rational endpoints of each $I_j$. The strict signs are
\begin{center}
\begin{tabular}{@{}lcccc@{}}\toprule
Shooting interval & lower phase 1 & lower phase 2 & upper phase 1 & upper phase 2\\\midrule
$I_1$ & $-$ & $+$ & $+$ & $-$\\
$I_2$ & $-$ & $+$ & $+$ & $-$\\
$I_3$ & $+$ & $+$ & $+$ & $-$\\\bottomrule
\end{tabular}
\end{center}
Here each entry holds at both endpoints, not just at one sampled orbit. Since $x(t;y)$ is strictly increasing in $y$ at every fixed time, each strict phase-level sign remains unchanged for all $y\in I_j$. No grazing is possible anywhere in these intervals, and no crossing can be lost at $t=0$ or $2\pi$, where the states are strictly inside the middle region. Thus all three intervals retain their certified transverse itineraries. The crossing times, and therefore $D$, are real analytic on a neighborhood of every $I_j$.

It follows that the zero set of $D$ within each $I_j$ is discrete: an analytic function with two nonzero endpoint values cannot vanish identically. We therefore obtain \emph{three distinct isolated fixed points}, not just three zeros belonging to a possible continuum of fixed points.

\subsection{Symmetry and isolation at the rational witness}
The forcing possesses the half-period reversal symmetry
\begin{equation}\label{eq:sym}
\mathcal Sx(t)=-x(t+\pi).
\end{equation}
The oddness of $g$ implies that $\mathcal S$ carries any periodic solution to another periodic solution. There is exactly one solution fixed by $\mathcal S$: the antisymmetric (central) solution.

To see existence and uniqueness directly, define
\[
Q(y)=x(\pi;0,y)+y.
\]
Scalar order makes $Q$ strictly increasing. Since $g(x)=9x+O(1)$ at both ends, $x(\pi;0,y)=e^{-9\pi}y+O(1)$ as $y\to\pm\infty$, so $Q$ has exactly one zero $y_c$. For that initial value, uniqueness and symmetry give $x_c(t+\pi)=-x_c(t)$ and therefore $2\pi$-periodicity. Our interval certificate gives
\begin{equation}\label{eq:Qbound}
Q(-0.571)>0.01911577,
\end{equation}
so $y_c<-0.571$, below all three intervals in \eqref{eq:intervals}. By scalar order, the corresponding three periodic solutions lie strictly \emph{above} $x_c$ at every common phase. Their three reflections under $\mathcal S$ lie strictly \emph{below} $x_c$, and they are distinct from one another because the flow preserves strict initial-value order. Hence we have $3+1+3$ distinct periodic orbits.

It remains to check that the central orbit itself is isolated; otherwise it would not count as a limit cycle. The interval computation additionally proves
\begin{equation}\label{eq:centerbracket}
-0.579022471<y_c<-0.579022470
\end{equation}
and bounds the lower-zone residence $s$ during $[0,\pi]$ by
\begin{equation}\label{eq:sbound}
s<0.448589954<\frac\pi7.
\end{equation}
Half-period antisymmetry makes its total exterior residence $2s$. The exact multiplier formula \eqref{eq:multiplier} therefore yields
\begin{equation}\label{eq:centerhyper}
\log P'(y_c)=\frac32(2\pi-2s)-9(2s)=3\pi-21s>0.
\end{equation}
Thus $P'(y_c)>1$; the central fixed point is hyperbolic and isolated. The symmetry map preserves isolation of each noncentral fixed point, completing the proof of Proposition~\ref{prop:witness}.

% Quantitative consequences of the strict seven-cycle interval certificate.
\subsection{Parameter-uniform isolation and stability}\label{sec:consequences}
Write $\theta=(A,B,\mu)$ and
$\theta_0=(9,3/2,181363/100000)$.
Let $x_\theta(t;y)$ solve the initial-value problem with
$x_\theta(0;y)=y$, and put
$D_\theta(y)=x_\theta(2\pi;y)-y$.

Theorem~\ref{thm:main} already states the uniform result.
Its proof reduces to three checks: preservation of the
six Poincar\'e displacement signs, absence of grazing on
the three shooting intervals, and the hyperbolicity of
the central antisymmetric orbit.

\begin{proof}[Completion of the proof of Theorem~\ref{thm:main}]
We use the three ingredients described above, with all
inequalities uniform over the cube~\eqref{eq:param-cube}.

\emph{1. Uniform displacement control.}
The reference vector field points strictly inward on the
boundaries $x=\pm2$, since $15/2-\mu_0>0$.
Consequently $|x_{\theta_0}(t;y)|\le2$ for all
$|y|\le2$ and $t\ge0$.
For $\theta$ in \eqref{eq:param-cube}, all secant slopes
of the perturbed scalar field $f_\theta(\cdot,t)$ are
at most $B<2$, whereas on each reference orbit
\[
 |f_\theta(x_{\theta_0},t)-f_{\theta_0}(x_{\theta_0},t)|
 \le|\Delta A|(|x_{\theta_0}|+|\sat(x_{\theta_0})|)
 +|\Delta B||\sat(x_{\theta_0})|+|\Delta\mu|
 <5\cdot10^{-12}.
\]
The scalar one-sided Gronwall inequality for solutions
started from the same initial value thus gives
\begin{equation}\label{eq:uniform-parameter}
 |x_\theta(t;y)-x_{\theta_0}(t;y)|
 \le5\cdot10^{-12}\frac{e^{2t}-1}{2}
 <1.5\cdot10^{-6},\quad
 0\le t\le2\pi,\quad |y|\le2.
\end{equation}
Indeed, $\pi<13/4$ and $e<11/4$ imply
$e^{4\pi}<e^{13}<(11/4)^{13}<600000$.
The smallest absolute certified displacement endpoint
margin in Table~\ref{tab:signs} exceeds
$4.96788\cdot10^{-6}$, so all six signs persist.
The intermediate value theorem therefore gives at least
one periodic solution in each of $I_1,I_2,I_3$.
At the central parameter triple, their existence and
the strict switching pattern are established by
Proposition~\ref{prop:witness}.

\emph{2. Explicit grazing barriers and isolation.}
At a tangency to $\ell\in\{-1,1\}$ the phase
$t_*$ must satisfy $\mu\sin t_*+B\ell=0$.
For the parameters under consideration $\mu>B>0$.
There are exactly four such phases in a period,
and each lies within $10^{-10}$ of its reference
phase: the perturbation of the phase equation
at the old root is less than $2\cdot10^{-12}$,
whereas the modulus of its time derivative
is greater than $0.8$ in the relevant
$10^{-10}$ neighborhood.

The original exact-flow directed-interval calculation
for the four phases and the \emph{six} rational shooting
endpoints, as now checked by
\texttt{proof/certify\_parameter\_persistence.py},
establishes the stronger quantitative bound
\begin{equation}\label{eq:grazingmargin}
 \bigl|x_{\theta_0}(t_*^{\,0};y)-\ell\bigr|
 >3\cdot10^{-4}
\end{equation}
at every relevant phase $t_*^{\,0}$, endpoint $y$,
and associated threshold $\ell$.
This inequality concerns absolute distances; their
signs are those in the grazing-barrier table of
Section~\ref{sec:proof}.
The scalar velocity is bounded by $40$ on the
trajectories under consideration, since
$|x_\theta|<2.000002$. Combining this with
\eqref{eq:uniform-parameter} gives
\[
 \bigl|x_\theta(t_*^\theta;y)-
              x_{\theta_0}(t_*^{\,0};y)\bigr|
 <1.5\cdot10^{-6}+40\cdot10^{-10}
 <1.505\cdot10^{-6}.
\]
Thus every strict barrier sign persists at
each endpoint. Since $y\mapsto x_\theta(t;y)$
is strictly increasing, none of the trajectories
with $y\in I_j$ can graze a switching level.
The shooting trajectories also stay uniformly
away from switching levels at $t=0,2\pi$;
there can be no crossing loss at the temporal
endpoints. Continuing parameters inside the
connected cube preserves the transverse
crossing itinerary on each $I_j$. Hence
$D_\theta$ is real analytic in $y$ on a
neighborhood of every $I_j$.
Opposite endpoint signs prevent the analytic
function from being identically zero.
Its zeros on $I_j$ are finite and isolated.

The three endpoint sign patterns are
$(+,-),(-,+),(+,-)$, respectively.
A positive-to-negative zero of a scalar
return displacement is locally attracting,
and a negative-to-positive zero is locally
repelling: the strictly increasing Poincar\'e
map moves nearby iterates monotonically
toward, respectively away from, the fixed point.
Therefore at least two attracting and
one repelling isolated periodic solutions
lie above the central antisymmetric solution.

\emph{3. Quantitative persistence of the central repeller.}
For all $\theta$ the equation has a unique
antisymmetric periodic orbit. Its initial
value $y_c(\theta)$ is the unique root of
$Q_\theta(y)=x_\theta(\pi;y)+y$, a strictly
increasing function with incremental slope
at least one. The base certificate encloses
$y_c(\theta_0)$ between
$-0.579022471$ and $-0.579022470$.
Using the improved bound $B<1.51$ up to
$t=\pi$, we have
\[
 |x_\theta(\pi;y)-x_{\theta_0}(\pi;y)|
 \le5\cdot10^{-12}\frac{e^{1.51\pi}-1}{1.51}
 <5\cdot10^{-12}\frac{158-1}{1.5}
 <6\cdot10^{-10},
\]
where $1.51\pi<5$ and $e^5<(11/4)^5<158$.
Consequently
$|y_c(\theta)-y_c(\theta_0)|<6\cdot10^{-10}$.
For $0\le t\le6/5$, the same estimate and
dependence on initial data yield
\[
 |x_\theta(t;y_c(\theta))-
        x_{\theta_0}(t;y_c(\theta_0))|
 < 5\cdot10^{-9},
\]
using $e^{1.51(6/5)}<e^2<(11/4)^2$.

The strict reference center shooting brackets,
together with monotonicity of downward and
upward crossing times in the shooting initial
value, place the two lower-zone events in
\[
 0.63458137<t_1^0<0.63458138,\qquad
 1.08317132<t_2^0<1.08317133.
\]
The threshold velocities $\mu_0\sin t-B_0$
are less than $-0.4$ on the first box and
greater than $0.09$ on the second.
For $|t-t_j^0|\le10^{-6}$, the absolute
change in $\dot x$ is less than
$410\cdot10^{-6}<0.001$, because
$|\dot x|<40$, $|f_x|<10$, and $\mu<2$.
Hence the reference level gaps at
$t_j^0\pm10^{-6}$ have absolute value
at least $0.089\cdot10^{-6}>5\cdot10^{-9}$.
Both crossing signs persist for the perturbed
central orbit, and its two crossing times
differ from the reference ones by less
than $10^{-6}$. The absence of further
lower-zone crossings follows from the
single-positive/single-negative phase arcs
of the threshold velocity (Lemma~\ref{lem:cross});
the upper level cannot be crossed and
recrossed on $[0,\pi]$, where its vector field
points upward. Thus its lower-zone residence
$s_\theta$ in the half-period obeys
\[
 s_\theta<(1.08317133-0.63458137)+2\cdot10^{-6}
 =0.44859196<0.448592.
\]
By antisymmetry the total exterior residence
in one period is $2s_\theta$, and the exact
Floquet exponent is
\begin{align*}
 \log P_\theta'(y_c(\theta))
 &=2\pi B-2(A+B)s_\theta\\
 &>2(3.14159)(1.5-10^{-12})
     -2(10.5+2\cdot10^{-12})(0.448592)
 >0.0043>0.
\end{align*}
The central orbit is therefore hyperbolic
repelling and isolated throughout the cube.

Finally, half-period reflection
$x(t)\mapsto-x(t+\pi)$ preserves each local
stability type, and provides three distinct
isolated counterparts below the central
orbit. This gives at least four attracting
and three repelling periodic orbits.
\end{proof}

\begin{corollary}[Critical residence-time profiles forced by the cycles]
\label{cor:critical}
For every parameter triple in \eqref{eq:param-cube},
the mean-parametrized biased response has at least
three distinct positive-mean stationary profiles
satisfying
\[
   \Lambda_\xi(\xi,\mu)=0,\qquad
   T(\xi,\mu)=\rho=\frac{2\pi A}{A+B}.
\]
\end{corollary}
\begin{proof}
Theorem~\ref{thm:main} gives three ordered
periodic orbits above the antisymmetric central
solution, with means $0<\xi_1<\xi_2<\xi_3$.
The central mean is zero and the biased
response vanishes at all four means.
Rolle's theorem applied to $(0,\xi_1)$,
$(\xi_1,\xi_2)$, and $(\xi_2,\xi_3)$
gives three distinct stationary profiles.
The identity \eqref{eq:key}, including the
continuous extension at an isolated grazing,
gives $T=\rho$ at all three.
Neither nondegeneracy of these points nor
an exact global count of folds is asserted.
\end{proof}

The supplemental script
\texttt{proof/certify\_parameter\_persistence.py}
verifies the directed interval margins in
\eqref{eq:grazingmargin}, the central event brackets,
and the rational inequalities used here. It reuses
the original directed-\texttt{Decimal} arithmetic;
it is not an independent Arb/FLINT implementation.


\section{Discussion and open questions}
The proof does not require numerically counting the turning points in Figure~\ref{fig:response}; those estimates were only a guide for locating the parameters. The numerical response curve exhibits an additional reversal near lower-level grazing. Proposition~\ref{prop:grazing} proves the local residence-time singularity that makes this geometry possible, but it does not prove the number of global response extrema. The reason we can \emph{prove} their existence is independent and simpler: three verified Poincar\'e sign changes and the half-period reflection symmetry.

Theorem~\ref{thm:main} disproves the conjectured five-cycle
bound on an explicitly defined open parameter cube.
The persistence proof certifies the switching structure
and the central multiplier uniformly over the cube; the
stability classification of noncentral cycles relies on
Poincar\'e displacement signs and order, and does not
establish their hyperbolicity. In particular the theorem
does not determine the sharp global upper bound, nor
the exact number of folds of the periodic response.
The general upper bound from~\cite{HuYuBound} remains unchanged.

The computations deserve independent scrutiny before journal submission. The current sign certificate relies on a standard-conforming Python \texttt{Decimal} implementation with correctly directed rounding; it is not a formally verified proof kernel. A second implementation based on Arb/FLINT or exact rational intervals would be a valuable independent audit. We separate this software-level qualification from the mathematical implication of the displayed strict bounds.

\section*{Code and data availability}
The public source repository contains this manuscript, the exact-flow interval checker, the exploratory response data, and independent high-precision and ODE checks. From the project directory run
\begin{verbatim}
python3 proof/certify_seven.py
python3 proof/certify_parameter_persistence.py
python3 -m unittest discover -s tests -v
\end{verbatim}
The certificate requires only the Python standard library. The file \texttt{response\_data.csv}, stored in \texttt{paper/figures/}, is a numerical plot sample, not a proof input. Neither the plot nor rounded roots are used for the existence theorem.

\appendix
\section{How the interval sign certificate works}\label{app:certificate}
We record the computational contract so that the argument is verifiable without trusting the location of the plotted roots.

\paragraph{Rational input and function enclosures.}
Every shooting initial value and all parameter coefficients are exact rationals. We enclose $\pi$ using Machin's identity
\[
\pi=16\arctan(1/5)-4\arctan(1/239)
\]
with alternating rational arctangent series and first-omitted-term bounds. The values of $e^z$, $\sin z$, and $\cos z$ appearing in \eqref{eq:flight} are evaluated by interval Taylor sums with outward rounding. In all calls, $|z|<10$; truncation at order at least $160$ admits the uniform error enclosure
\begin{equation}\label{eq:tail}
\bigl|\operatorname{remainder}\bigr|\le\frac{3^{10}10^{161}}{161!}<10^{-95}.
\end{equation}
Products use all four endpoint combinations, and interval reciprocals are allowed only when their denominator excludes zero. The actual decimal arithmetic is performed with 95 digits and separate downward and upward rounding modes.

\paragraph{Certified event times.}
The exploratory solver supplies decimal \emph{centers} for event times, not trusted root values. For each event, the checker tests an interval $[t_-,t_+]$ with rational endpoints. It evaluates the exact flight \eqref{eq:flight} at both endpoints and requires strict opposite signs relative to the relevant level $\ell=\pm1$; the interval crossing speed has a prescribed nonzero sign. The velocity shape identity \eqref{eq:shape} prevents a hidden earlier encounter. The four possible grazing phases \eqref{eq:grazingphases} are separately bracketed, and the signs of $x(t;y)-\ell$ at each phase are checked at both ends of every shooting interval. Scalar order then excludes grazing throughout the intervals. For two-crossing cases, the initial minimum is separately enclosed above $-1$. The exact state at the enclosed event is the known switching level $\ell$, so the next flight begins from the interval event time and that \emph{exact} state. Interval propagation continues until $2\pi$, producing Table~\ref{tab:signs}.

\paragraph{The central orbit.}
The same half-period algorithm evaluates $Q(y)=x(\pi;0,y)+y$ at the rational endpoints of \eqref{eq:centerbracket}, with opposite certified signs. After checking the two transverse lower crossings, it encloses their time separation at the lower endpoint of the center bracket. Monotonic order of the first descending and second ascending crossing times shows that this is an upper bound for the central lower-zone residence $s$. The strict rational comparison $s<\pi/7$ gives \eqref{eq:centerhyper}.

\paragraph{Logical dependencies.}
The executable interval certificates are used for four groups of
strict inequalities: (i) the six displacement signs,
(ii) the crossing and grazing barriers on the three shooting
intervals, (iii) the central shooting and residence-time
bounds, and (iv) the margins in the quantitative parameter
continuation. The sign-change, order, symmetry, analyticity,
Gronwall, and Floquet arguments are mathematical deductions
from those inequalities. The rounded numerical multipliers
in Table~\ref{tab:roots} are not proof inputs.

\begin{thebibliography}{9}
\bibitem{Bravo2025}
J.~L. Bravo, V.~Carmona, M.~Fern\'andez and I.~Ojeda,
\emph{Upper bound of the number of limit cycles for symmetric scalar piecewise linear differential equations with three zones}, arXiv:2507.06705 (2025).
\bibitem{BravoTrinidad2026}
J.~L. Bravo and R.~Trinidad-Forte,
\emph{Finitude of limit cycles of linear piecewise ODEs in the cylinder}, arXiv:2605.05805 (2026).
\bibitem{HuYuBound}
C.~Hu and T.~Yu,
\emph{Sparse Toric Flight Encoding for Constant-Slope Scalar Periodic Piecewise-Linear Equations, with an Explicit Three-Zone Limit-Cycle Bound}, manuscript (2026).
\end{thebibliography}
\end{document}
